%% kapa: exam=1 round=30 year=2019 period=2 form=A subject=appraisal_law subject_ko=감정평가관계법규 seq=17 number=17
%% source: Q-Net 감정평가사 기출문제 (cst00309) artlSeq=5205205 file=extracted/1008/5205205/2019년 제30회 감정평가사 1차시험 문제지/2019년 제30회 감정평가사 1차시험 문제지/2019년도 제30회 감정평가사 1차 2교시 A형.hwp
%% conversion: kordoc 4.24.51 -> tools/gen1.py (auto); review=proofread-sonnet; flags=none
\begin{question}{17}
부동산 가격공시에 관한 법령상 개별공시지가에 관한 설명으로 옳지 않은 것은?
\begin{choices}
  \item 개별공시지가에 이의가 있는 자는 그 결정ㆍ공시일부터 30일 이내에 서면으로 시장ㆍ군수 또는 구청장에게 이의를 신청할 수 있다.
  \item 시장ㆍ군수 또는 구청장은 개별공시지가에 표준지 선정의 착오가 있음을 발견한 때에는 지체 없이 이를 정정하여야 한다.
  \item 관계 공무원이 표준지가격의 조사ㆍ평가를 위하여 택지에 출입하고자 할 때에는 시장ㆍ군수 또는 구청장의 허가를 받아 출입할 날의 3일 전에 그 점유자에게 일시와 장소를 통지하여야 한다.
  \item 일출 전ㆍ일몰 후에는 그 토지의 점유자의 승인 없이 택지 또는 담장이나 울타리로 둘러싸인 타인의 토지에 출입할 수 없다.
  \item 개별공시지가의 결정ㆍ공시에 소요되는 비용은 그 비용의 50퍼센트 이내에서 국고에서 보조할 수 있다.
\end{choices}
\end{question}
\answer{3}
